\subsection{FPTRANSA ULP Error Model}
For an approximate transcendental instruction, $x'$ is the source after \texttt{FSTATUS.DAZ}, $y=f(x')$ is the
exact real reference value named by its accuracy contract, and $a$ is the approximate result before
\texttt{FSTATUS.FTZ}. For a format with precision $p$ and minimum normal exponent $e_{\min}$, define

\[
  \operatorname{ulp}_F(y) =
  \begin{cases}
    2^{\max(\lfloor\log_2 |y|\rfloor-p+1,\ e_{\min}-p+1)}, & y\ne0,\\
    2^{e_{\min}-p+1}, & y=0.
  \end{cases}
\]

The measured error is $|a-y|/\operatorname{ulp}_F(y)$. The bound applies only to finite reference values within
the selected format's finite range; NaN, infinity, and overflow follow the instruction's special-value rules.
Every present contract guarantees at most 4 ULP for both S and D, corresponding to at least 21 and 50 worst-case
relative precision bits respectively for nonzero normal results. CPUID may advertise a smaller certified bound.
The instruction does not compute or return its actual error.

FPTRANSA applies DAZ before forming the exact reference input and measures its advertised ULP error before FTZ.
It ignores \texttt{FSTATUS.RM} and formats its implementation-defined approximation with nearest-even rounding.
The approximation is deterministic for the same implementation, input, format, and relevant FSTATUS mode, but need
not be bit-identical across implementations. DN and FTZ retain their ordinary meanings. FPTRANSA may accrue NV, DZ,
OF, and UF as its instruction contract requires, but it leaves the accrued NX flag unchanged and never takes an NX
trap. Strict language-library operations may use FPTRANSA only through an explicit approximate intrinsic or a
compiler mode that permits approximate math; strict \texttt{libm} lowering is not valid.
